a) Jede nilpotente Matrix
\mathl{\begin{pmatrix}  a   &   b   \\  c  & d \end{pmatrix}}{} lässt sich durch den linearen Weg
\maabbeledisp {} {[0,1]} {N
} { t } {t \begin{pmatrix}  a   &   b   \\  c  & d \end{pmatrix}
} {,}
innerhalb der nilpotenten Matrizen mit der Nullmatrix verbinden. Daher ist $N$ wegzusammenhängend und damit auch zusammenhängend.

b) Eine
$2 \times 2$-\definitionsverweis {Matrix}{}{}
ist genau dann nilpotent, wenn sowohl die Spur als auch die Determinante $0$ sind. Die Menge $N$ der nilpotenten Matrizen kann also als

\mathdisp {\left\{  (a,b,c,d)  \mid   a+d = 0 \text{ und } ad-bc = 0         \right\}} {  }

aufgefasst werden. Wir betrachten die Abbildung
\maabbeledisp {} {\R^4} {\R^2
} {(a,b,c,d)} { (a+d, ad-bc)
} {.}
Deren Jacobi-Matrix ist

\mathdisp {\begin{pmatrix}  1 &  0 &  0 &  1 \\  d  &  -c &  -b &  a  \end{pmatrix}} { . }

Diese Abbildung ist im Nullpunkt nicht
\definitionsverweis {regulär}{}{,}
aber in jedem anderen Punkt der Faser $N$. Wenn nämlich
\mathl{P \in M}{} ist, so folgt wegen

\mavergleichskettedisp
{\vergleichskette
{a^2
}
{ =} {bc
}
{ } {
}
{ } {
}
{ } {
}
}
{}{}{}
aus

\mavergleichskettedisp
{\vergleichskette
{b
}
{ =} { c
}
{ =} { 0
}
{ } {
}
{ } {
}
}
{}{}{}
sofort

\mavergleichskettedisp
{\vergleichskette
{a
}
{ =} {b
}
{ =} { c
}
{ =} { d
}
{ =} { 0
}
}
{}{}{.}
Die Jacobi-Matrix hat also in den Punkten aus $M$ den maximalen Rang. Mit

\mavergleichskettedisp
{\vergleichskette
{G
}
{ =} {\R^4 \setminus \{0\}
}
{ } {
}
{ } {
}
{ } {
}
}
{}{}{}
kann man $M$ als die Faser der eingeschränkten Abbildung auffassen, die überall auf der Faser regulär ist. Daher ist $M$ nach
[[Satz über implizite Abbildungen/Faser ist Mannigfaltigkeit/Fakt|Fakt]]
eine abgeschlossene Untermannigfaltigkeit von $G$.

c) Nach
[[Satz über implizite Abbildungen/Faser ist Mannigfaltigkeit/Fakt|Fakt]]
ist die Dimension von $M$ gleich

\mavergleichskette
{\vergleichskette
{ 4-2
}
{ = }{ 2
}
{  }{
}
{    }{
}
{   }{
}
}
{}{}{.}

d) Wir schreiben

\mavergleichskettedisp
{\vergleichskette
{M_+
}
{ =} { \left\{ (a,b,c,d) \in M  \mid   b>c        \right\}
}
{ } {
}
{ } {
}
{ } {
}
}
{}{}{}
und

\mavergleichskettedisp
{\vergleichskette
{M_-
}
{ =} { \left\{ (a,b,c,d) \in M  \mid   b<c        \right\}
}
{ } {
}
{ } {
}
{ } {
}
}
{}{}{,}
beides sind
\zusatzklammer {als Durchschnitt von $M$ mit der durch
\mathl{b>c}{} gegebenen offenen Menge des $\R^4$} {} {}
offene Mengen in $M$. Die Matrizen
\mathkor {} {\begin{pmatrix}  0  &   1  \\  0 &  0 \end{pmatrix}} {und} {\begin{pmatrix}  0  &   0  \\  1 &  0 \end{pmatrix}} {}
zeigen, dass sie nicht leer sind. Ferner überdecken sie ganz $M$. Bei

\mavergleichskette
{\vergleichskette
{b
}
{ = }{ c
}
{  }{
}
{    }{
}
{   }{
}
}
{}{}{}
folgt nämlich wegen

\mavergleichskettedisp
{\vergleichskette
{ 0
}
{ =} { ad-bc
}
{ =} { -a^2 -b^2
}
{ } {
}
{ } {
}
}
{}{}{}
direkt

\mavergleichskette
{\vergleichskette
{a
}
{ = }{b
}
{ = }{ c
}
{ =   }{ d
}
{ =  }{ 0
}
}
{}{}{,}
und der Punkt gehört nicht zu $M$. Es liegt also eine Überdeckung mit zwei nichtleeren disjunkten offenen Mengen vor, daher ist $M$ nicht zusammenhängend.

e) Wir arbeiten mit der Abbildung
\maabbeledisp {\psi} {\R^2 \setminus \{(0,0)\}} { M_+
} {(a,u)} { (a, u + \sqrt{a^2+u^2}, u - \sqrt{a^2 +u^2},-a)
} {.}
Wegen

\mavergleichskettealign
{\vergleichskettealign
{ad-bc
}
{ =} { -a^2 - \left(  u + \sqrt{a^2+u^2}  \right) \left( u - \sqrt{a^2+u^2}  \right)
}
{ =} { -a^2 - \left(  u^2 - \left(  a^2+u^2 \right)   \right)
}
{ =} { 0
}
{ } {
}
}
{}
{}{}
ist die Determinantenbedingung erfüllt, und wegen

\mavergleichskettedisp
{\vergleichskette
{b-c
}
{ =} { u + \sqrt{a^2+u^2} - \left(  u- \sqrt{a^2 +u^2})  \right)
}
{ =} { 2 \sqrt{a^2+u^2}
}
{ >} { 0
}
{ } {
}
}
{}{}{}
gehört das Bild zu $M_+$. Die Abbildung
\maabbeledisp {\varphi} {M_+} { \R^2 \setminus \{(0,0)\}
} {(a,b,c,d)} { \left(  a  , \,   { \frac{  b+c }{  2 } }                   \right)
} {,}
ist invers zu der gegebenen Abbildung. Dabei ist

\mavergleichskettedisp
{\vergleichskette
{ \varphi \circ \psi
}
{ =} {
\operatorname{Id}_{ \R^2 \setminus \{(0,0)\}  }
}
{ } {
}
{ } {
}
{ } {
}
}
{}{}{}
klar. Die andere Identität ergibt sich aus

\mavergleichskettealign
{\vergleichskettealign
{ { \frac{  b+c }{  2 } } + \sqrt{a^2 + \left( { \frac{  b+c }{  2 } }  \right)^2 }
}
{ =} { { \frac{  b+c }{  2 } } + { \frac{   1 }{  2 } }  \sqrt{4 a^2 + b^2 +c^2 +2bc }
}
{ =} { { \frac{  b+c }{  2 } } + { \frac{   1 }{  2 } }  \sqrt{ b^2 +c^2 -2bc }
}
{ =} { { \frac{  b+c }{  2 } } + { \frac{  b-c }{  2 } }
}
{ =} {b
}
}
{}
{}{}
und

\mavergleichskettedisp
{\vergleichskette
{ { \frac{  b+c }{  2 } } - \sqrt{a^2 + \left(   { \frac{  b+c }{  2 } }   \right)^2 }
}
{ =} { { \frac{  b+c }{  2 } } - { \frac{  b-c }{  2 } }
}
{ =} { c
}
{ } {}
{ } {}
}
{}{}{.}
Beide Abbildungen sind stetig, daher liegt eine
\definitionsverweis {Homöomorphie}{}{}
vor.
Für
\mathl{M_-}{} vertauscht man die Rollen von
\mathkor {} {b} {und} {c} {.}

 [[Kategorie:Latexseite]]
